% LaTeX .tex example for the proceedings of
% COBEM 2017 - 24th International Congress of Mechanical Engineering
% December, 3-8, 2017, Curitiba, PR, Brazil
%
% Based on the template of the proceedings of COBEM2015 
% MODELO COBEM 2017

\documentclass[draft]{svjour3}

\usepackage{engsymbols}
\usepackage{magref}
\usepackage[utf8]{inputenc}
\usepackage{nicefrac}

\newcommand{\bmax}{\ensuremath{{B}\ped{max}}}
\newcommand{\bmin}{\ensuremath{{B}\ped{min}}}
\newcommand{\mrate}{\rate{m}}
\newcommand{\wvalven}{\w\ped{valve,n}}
\newcommand{\wrelayn}{\w\ped{relay,n}}

\usepackage{mathptmx}
\journalname{Journal of the Brazilian Society of Mechanical Sciences and Engineering}

\usepackage{graphicx}
\usepackage{subfig}
\usepackage{xcolor}
\usepackage{nomencl}
\usepackage{hyperref}

\begin{document}

\title{Numerical analysis of the influence of magnetic field waveforms on the performance of active magnetic regenerators}

\author{
Fábio P. Fortkamp \and
Gusttav B. Lang \and
Jaime A. Lozano \and
Jader R. Barbosa Jr.
}


\institute{
Fábio P. Fortkamp 
\at
POLO --- Research Laboratories for Emerging Technologies in Cooling and Thermophysics, Department of Mechanical Engineering, Federal University of Santa Catarina, Florianópolis, SC, 88040-900, Brazil, \email{fabio@polo.ufsc.br}
\and
Gusttav B. Lang \and
Jaime A. Lozano \and
Jader R. Barbosa Jr. \at POLO --- Research Laboratories for Emerging Technologies in Cooling and Thermophysics, Department of Mechanical Engineering, Federal University of Santa Catarina, Florianópolis, SC, 88040-900, Brazil}


\date{}


\maketitle

\begin{abstract}


\keywords{ magnetic refrigeration \and active magnetic regenerator \and numerical modeling \and magnetocaloric effect}
\end{abstract}

\section{Introduction}
\label{sec:introduction}

Mechanical vapor compression has been the dominant cooling technology for the past century \cite{bib:nik} but, despite its dominance, it still faces an increasing number of challenges related to the its environmental footprint. The large entropy generation associated with vapor compression processes makes this techonology constrained by low exergy efficiencies, in particular for compact systems \cite{KITANOVSKI2015288}. In addition, the phase-out of refrigerants \textcolor{black}{with} ozone depleting \textcolor{black}{and global warming potentials has} introduced the use of flammable substances, with their own set of risks for consumer applications \cite{bib:iir-flammable} and technologial challenges \cite{bib:lionte18-adapt}.

The term \emph{caloric cooling} is \textcolor{black}{generally} applied to alternative cooling techonologies in which the cooling effect is generated by the response of a solid material to a physical \textcolor{black}{stimulus}. The absence of hazardous gases, compressors \textcolor{black}{and throttling devices} makes these technologies potentially more efficient, \textcolor{black}{scalable}, silent and environmentally friendly. Magnetic refrigeration (MR) is the most advanced of all caloric technologies \cite{KITANOVSKI2015288}; in particular, with the classes of magnetic refrigerants currently available, magnetocaloric refrigeration seems promising for applications in small temperature span (around \SI{20}{\kelvin}), such as winecoolers \textcolor{black}{and air conditioners} \cite{KITANOVSKI2015288}.

\nomenclature[vm]{MR}{Magnetic (or Magnetocaloric) Refrigeration}

In magnetic refrigeration, a \emph{magnetocaloric material} (MCM) \textcolor{black}{is subjected to a cyclical change of the applied} magnetic field, and its temperature \textcolor{black}{changes as a result} of the \emph{magnetocaloric effect} (MCE). The magnitude of the MCE  \textcolor{black}{depends} on  parameters \textcolor{black}{related to} characteristics of the material, magnetic field variation and temperature, and its maximal at the transition Curie temperature of the material. For a thorough introduction to the magnetocaloric effect and magnetocaloric materials, the reader is referred to \cite{bib:smith-magneto}.

\nomenclature[vm]{MCM}{Magnetocaloric Material}

For operating conditions typical of a low-temperature span magnetic refrigerator, the MCE is typically \textcolor{black}{on} the order of \num{2}-\SI{5}{\kelvin}. \textcolor{black}{T}o amplify this \textcolor{black}{temperature change}, heat regeneration is usually employed \cite{bib:kitanovski}. Active magnetic regenerators (AMR) are thermal devices where the magnetocaloric material constitutes a solid matrix through which flows an aqueous heat transfer fluid, and is cyclically magnetized and demagnetized to activate the magnetocaloric effect.  \textcolor{black}{A regenerator is essentially a cascade of infinitesimal ``layers'' of MCM that are activated simultaneously} to build up a temperature profile \textcolor{black}{along its length}. These layers can be made of the same material, yielding an homogeneous regenerator; however, given the  dependence  of the MCE with temperature, it is an interesting strategy to build \emph{multilayer} regenerators, where each layer is composed of a different material \textcolor{black}{that} will work around its \textcolor{black}{own} Curie temperature, maximizing the magnetocaloric effect of each portion.

\nomenclature[vm]{AMR}{Active Magnetic Regenerator}

The five main subsystems in an AMR-based magnetic refrigerator can be identified as follows \cite{bib:jaime,bib:trevizoli16_pump}:

\begin{enumerate}
\item magnetic circuit; \label{item:1}
\item active magnetic regenerator; \label{item:2}
\item flow management system; \label{item:3}
\item control system; \label{item:4}
\item transmission system; \label{item:5}
\item heat exchangers and cabinet. \label{item:6a}
\end{enumerate}

The magnetic circuit generates oscillating magnetic fields that activate the MCE \textcolor{black}{in} the active magnetic regenerators. As previously explained, the regenerators are composed of magnetocaloric porous media, and the flow management system provides the alternating fluid flow through the beds. The control system \textcolor{black}{serves} the purpose of synchronizing the fluid flow regime with the magnetic field variations. Finally, the heat exchangers are responsible for transferring the cooling effect from the AMR to the refrigerated cabinet and the rejected heat to the surroundings. 

Magnetic refrigerators can operate according to different thermodynamic cycles, but the stages are roughly as follows. The magnetic circuit magnetizes the MCM, increasing its temperature due to the MCE. During the so-called \emph{cold blow}, cold fluid from the cold source flows through the warm bed and absorbs its energy, releasing it to the external heat exchanger. On the return of the fluid, the solid is demagnetized and cooled down, and in the \emph{hot blow} the fluid relases energy to the bed, decreasing its temperature so it can absorb the thermal load at the cold source. The cycle then repeats; different cycles can be configured by different durations and synchronization steps between these two \emph{characteristic waveforms}:

\begin{enumerate}
\item The \emph{magnetic profile}, which describes the oscillating magnetic field over one regenerator;
\item The \emph{fluid flow profile}, which describes the time-variation of flow rate through one bed.
\end{enumerate}

Among the challenges cited in the literature for further improvement of magnetic refrigerators, two can be highlighted \cite{bib:eriksen16_activ}:

\begin{enumerate}
\item Better use of AMR and magnetic circuit models during the design process, to identify optimal geometric and operating parameter;
\item Better study of losses that are present in a magnetic cooling device.
\end{enumerate}

\textcolor{black}{In} other words, these challenges arise due to \emph{interactions} between different subsystems \textcolor{black}of a magnetic refrigerator. In these context, the focus on the aforementioned waveforms is a suitable approach for optimizing the performance of magnetic refrigerators, as these profiles represent interactions between systems. As it will be seen, mathematical models for AMRs take these waveforms as input, and they are generated by models for other components.

Fluid flow profiles are more easily investigated with experimental methods, since they can be altered with modifications in the fluid flow system. In particular, the \emph{duration} of blows have been extensively investigated \cite{bib:teyber17_exper,bib:nakashima18-influen-exp,FORTKAMP2018}, and the literature shows that the cooling capacity of magnetic refrigerators can be increased by focusing fluid flows during periods where the magnetic field is at its extrema values.

In this work, we focus on magnetic profiles, which are less studied in the literature. The studied waveforms are depicted in \autoref{fig:profiles}, identified by the transitions between states. The instantaneous profile represents step-like changes between high and low levels of the magnetic field, while the ramp profile has linear variations between these states. The rectified cosine profile is a common waveform found in the literature that has a continuous variation of the applied magnetic field.

\begin{figure}[!ht]
  \centering
  \includegraphics[width=8cm]{profiles_all}
  \caption{Magnetic profiles studied in this paper}
  \label{fig:profiles}
\end{figure}

The instantaneous magnetic profile is characteristic of the Brayton AMR cycle, where the magnetization and demagnetization steps must be done adiabatically, as shown in \autoref{fig:brayton}; compared to other common cycles such as the Ericsson, the Brayton cycle results in the highest cooling capacities for the same operating  parameters \cite{bib:kitanovski}, at the expense of lower coefficient of performance.


\begin{figure}[!ht]
  \centering
  
  \caption{Brayton AMR cycle}
  \label{fig:brayton}
\end{figure}

 In practice, the ramp profile is a more feasible realization with finite transition times, and its synchronization with the fluid flow profile can yield different performance levels \cite{TUSEK20111507}. In particular, the plateaus of magnetic field should be centered with the blow durations \cite{bib:bjoerk11_amr}.
\autoref{fig:sync} shows the syncronization schema between the magnetic and fluid flow profiles, showing the ramp magnetic profile as an example. 

\begin{figure}[!ht]
  \centering
  \includegraphics[width=8cm]{profiles_rm_and_flow_instantaneous}
  \caption{Ramp magnetic profile (solid lines) and instantaneous fluid flow profile (dashed lines)}
  \label{fig:sync}
\end{figure}

This paper uses the AMR model developed by \cite{bib:trevizoli16_perfor_model}, and the same authors had studied magnetic profiles waveforms in an earlier work \cite{bib:asme-mce}. However, their amplitudes was not varied, nor were they investigated with different geometries.

The present works combines the gaps from the above cited papers in investigating the performance of a magnetic refrigerator under different magnetic profile waveforms. These waveforms are mathematically modeled, and the cooling capacity and coefficient of performance are calculated based on the profiles parameters, while also investigated how the AMR geometry affect the performance in combinations with the magnetic profile. The fluid flow profile is assumed fixed in shape, although its parameters are also variied. The emulate constraints on an operating point of actual magnetic refrigeration devices, the temperature span is set fixed, and hence few comments are made on second-law efficiency.



\section{Materials and methods}
\label{sec:materials-methods}

In this work, we perform numerical simulations of a previously developed AMR model, varying the profile and geometric parameters. We also couple a model for calculating  the power comsumption of a proposed valve system, a topic understudied in the literature. With this, we can evaluate the influence of the profile on cooling capacity and coefficient of performance.

\subsection{AMR model}
\label{sec:amr-model}

\subsubsection{Governing equations}
\label{sec:governing-equations}

Simulations were performed using a one-dimensional AMR mathematical model \cite{bib:asme-mce,bib:trevizoli16_perfor_model}. Here we will show a summarized version focusing on the governing equations. This model solves momentum and energy balance equations for the solid and fluid phases, represented by indices 's' and 'f', respectively. The model geometry is depicted in \autoref{fig:amrmodel}, and assumes regenerators composed of monodisperse packed spheres with porosity $\varepsilon$.

\nomenclature[is]{s}{solid phase}
\nomenclature[if]{f}{fluid phase}

\begin{figure}[!ht]
  \centering
  \includegraphics[width=8cm]{reg3d}
  \caption{AMR model geometry}
  \label{fig:amrmodel}
\end{figure}

The momentum equation for the fluid domain is:

\begin{equation}
\label{eq:75}
  \frac{\rho\ped{f}}{\varepsilon} \diffp{V_z}{t}= -\diffp{P}{z} - \frac{\mu\ped{f}}{K} V_z - \frac{c\ped{E} \rho\ped{f}}{K^{1/2}} \left| V_z\right| V_z 
\end{equation}

\nomenclature[ak]{$K$}{permeability of the porous medium [\si{\meter\square}]}
\nomenclature[ac]{$c\ped{E}$}{Ergun constant of the porous medium}
\nomenclature[ar]{$\rho$}{density [\si{\kg\per\cubic\meter}]}
\nomenclature[be]{$\varepsilon$}{porosity}
\nomenclature[av]{$V$}{velocity [\si{\meter\per\second}]}
\nomenclature[ap]{$P$}{pressure [\si{\pascal}]}
\nomenclature[bu]{$\mu$}{dynamic viscosity [\si{\pascal\second}]}


\noindent and is solved for the time-dependent uniform fluid velocity through the bed. 

The energy equation for the fluid phase can be written as:

\begin{equation}
\label{eq:78}
\begin{split}
  \rho\ped{f} c_{p,\mathrm{f}} \left(\varepsilon \diffp{T\ped{f}}{t} + V_z\diffp{T\ped{f}}{z}\right) = & -\hheat\ped{sf} \beta \left(T\ped{f} - T\ped{s}\right) \\
& + \left|V_z\diffp{p}{z}\bigg|_f \right| \\
& + \varepsilon\left(k\ped{f}\ap{eff}  +\rho\ped{f} c_{p,\mathrm{f}} D\ped{ld}\right)\diffp[2]{T\ped{f}}{z} \\
& + \rate{q}\ped{csg}  
\end{split}
\end{equation}

\nomenclature[at]{$T$}{temperature [\si{\kelvin}]}
\nomenclature[at]{$t$}{time [\si{\second}]}
\nomenclature[az]{$z$}{axial coordinate [\si{\meter}]}

\nomenclature[aq]{$\rate{q}\ped{csg}$}{volumetric casing losses in the AMR model [\si{\watt\per\cubic\meter}]}
\nomenclature[ic]{csg}{regenerator casing}
\nomenclature[ah]{$\hheat$}{heat transfer coefficient [\si{\watt\per\meter\squared\per\kelvin}]}
\nomenclature[aa]{$A\ped{sf}$}{contact surface area between solid and fluid phases in the regenerator [\si{\meter^2}]}
\nomenclature[is]{sf}{relative to the heat transfer between solid and fluid phases in the regenerator}
\nomenclature[azb]{$\beta$}{surface area density [\si{\squared\meter\per\cubic\meter}]}
\nomenclature[ak]{$k\ap{eff}$}{effective thermal conductivity [\si{\watt\per\meter\per\kelvin}]}
\nomenclature[ad]{$D\ped{ld}$}{dispersion term in the AMR model [\si{\meter\square\per\second}]}
\nomenclature[af]{$f$}{cycle frequency [\si{\hertz}]}
\nomenclature[azf]{$\Phi$}{utilization}
\nomenclature[am]{$m$}{mass [\si{\kg}]}

\nomenclature[azt]{$\tau$}{time period [\si{\second}]}
\nomenclature[ac]{$c$}{specific heat [\si{\joule\per\kg\per\kelvin}]}


\noindent while the energy equation for the solid phase is written as:

\begin{equation}
\label{eq:88}
  \rho\ped{s} c\ped{s}(1 - \varepsilon)\diffp{T\ped{s}}{t} = \hheat\ped{sf}\beta(T\ped{f}-T\ped{s}) + (1-\varepsilon)k\ped{s}\ap{eff}\diffp[2]{T\ped{s}}{z}
\end{equation}

Initial and boundary conditions, solution methods, convergence analyses, and closure relations for the porous media terms are described in \cite{bib:trevizoli16_perfor_model}. This AMR model solves the above equations for one regenerator operating between given sources temperatures, during one full cycle (hot and cold blows and  magnetization and demagnetization periods), given specified operating conditions (to be discussed later).

The casing heat transfer term $\rate{q}\ped{csg}$ in \autoref{eq:78} is calculated solving the heat conduction in the regenerators casing \cite{bib:trevizoli16_perfor_model}, and can be neglected to simplify some analyses. More details are provided in \autoref{sec:results-discussions}.

\subsubsection{How the fluid flow profile is used}
\label{sec:how-fluid-flow}

The pressure gradient in \autoref{eq:75}is modeled as:

\begin{equation}
\label{eq:76}
-\diffp{P}{z} = \rho\ped{f} A\ped{t} g(t)
\end{equation}

\noindent where \(g(t)\) is a dimensionless function that expresses the mathematical waveform of the pressure gradient, and  \(A\ped{t}\) is its amplitude,  adjusted in a convergence loop; the mass flow rate calculated from the Darcy velocity from \autoref{eq:75} is compared with the input value of mass flow rate until these values converge. The waveform $g(t)$ represents the fluid flow profile used.

\nomenclature[ag]{$g(t)$}{dimensional waveform of the pressure gradient term in the fluid momentum equation}
\nomenclature[aA]{$A\ped{t}$}{amplitude of the pressure gradient waveform in the fluid momentum equation [\si{\meter\per\second\square}]}


\subsubsection{How the magnetic profile is used}
\label{sec:how-magnetic-profile}


The magnetic profile is modeled by a waveform of magnetic field strength $H$ applied perpendicular to the regenerators, as shown in \autoref{fig:amrmodel}. The magnetic field is assumed uniform throughout the beds. The applied field is corrected from demagnetization effects to yield the effective field inside the regenerators:

\begin{equation}
  \label{eq:6}
  H\ap{eff} = H - N\ped{D} M
\end{equation}

\nomenclature[am]{$M$}{magnetization field [\si{\ampere\per\meter}]}
\nomenclature[an]{$N\ped{D}$}{demagnetization tensor}

\noindent where $M$ is the magnetization field of the material, and $N\ped{D}$ is a demagnetization factor.

The magnetocaloric effect is implemented in the so-called discrete approach \cite{bib:nielsen11_review}; every time the magnetic field changes, based on the input magnetic profile, the solid temperature is calculated according to:

\begin{equation}
  \label{eq:89}
  T\ped{s}(t + \Delta{}t) = T\ped{s}(t) + \dtad \left(T\ped{s}(t),H\ap{eff}(t),H\ap{eff}\left(t + \Delta{}t\right)\right)
\end{equation}


\nomenclature[ct]{$\dtad$}{adiabatic temperature variation [\si{\kelvin}]}

\noindent where the adiabatic temperature variation $\dtad$ is calculated from \textcolor{black}{tabulated} experimental \textcolor{black}{data} for magnetocaloric materials as function of temperature and effective field. \cite{bib:trevizoli16_perfor_model}.

\subsubsection{Evaluation of solid and fluid properties}
\label{sec:eval-solid-fluid}

The fluid properties are considered constant in the momentum equation to decouple the solution \textcolor{black}{procedures to determine} the velocity and temperature fields. The properties are computed at the average temperature between the hot and cold sources temperatures, and are evaluated from interpolated tables exported from the EES software \cite{bib:klein13-ees}. In all simulations shown in this work, the heat transfer fluid is a mixture of water/ethylene-glycol with concentration \num{80}/\SI{20}{\percent} vol. For the energy equations, the fluid properties are also calculated from tabulated data, but the temperature dependence is considered.

n this chapter, all simulations use gadolinium or its allows as the magnetocaloric material. Gadolinium is a benchmark material with a Curie temperature of \SI{290}{\kelvin}. \textcolor{black}{For simplicity, the solid} density is assumed constant at $\rho\ped{s} = \SI{7900}{\kg\per\cubic\meter}$ and the solid \textcolor{black}{thermal} conductivity is set to $k\ped{s} = \SI{10.5}{\watt\per\meter\per\kelvin}$. The specific heat \textcolor{black}{capacity} of \ce{Gd} is calculated as a function of temperature and magnetic field based on experimental \textcolor{black}{data}, using a bi-linear interpolation scheme; more details on the experimental dataset \textcolor{black}{are available} in \cite{bib:trevizoli16_perfor_model}. 

In this work, both single- and multilayer regenerators are considered. For the multilayer simulations, alloys of gadolinium and yttrium are used in the form \ce{Gd_{1-x}Y_x}, where $x$ is the yttrium fraction.  \textcolor{black}{This fraction} decreases the Curie temperature  \textcolor{black}{of the} alloy relative to that of pure gadolinium. Due to the lack of experimental  \textcolor{black}{data on the magnetocaloric properties of} \ce{Gd_{1-x}Y_x} alloys at the time \textcolor{black}{this analysis was made}, a simpler approach was used \textcolor{black}{in which} the properties of alloys with {a} low yttrium fraction are \textcolor{black}{identical to those} of pure gadolinium, except for the Curie temperature \textcolor{black}{(which is shifted to lower values)}. 


\subsubsection{Performance metrics}
\label{sec:performance-metrics}

The AMR model considers only one bed and assumes $N\ped{r}$ identical beds experience the same cycle, multiplyind the performance metrics below by this factor

\nomenclature{an}{$N\ped{r}$}{number of regenerators}

The cooling capacity is calculated as \cite{bib:trevizoli16_perfor_model}:

\begin{equation}
\label{eq:99}
\qc = N\ped{r}\frac{1}{\tau} \int_{\tau\ped{HB}} \mrate\ped{f}(t) c_{p,\mathrm{f}} \left( T\ped{C} -T\ped{f,CE} \right) \diffd{t}
\end{equation}

\noindent where $\tau$ is the cycle period, and $\tau\ped{HB}$ is the duration of the hot blow, during which the cold end temperature is compared to the cold source temperature.

\nomenclature[vh]{HB}{hot blow}
\nomenclature[vc]{CE}{cold end}
\nomenclature[am]{$\mrate\ped{f}$}{mass flow rate [\si{\kg\per\second}]}
\nomenclature[vc]{C}{cold source}
\nomenclature[vc]{H}{hot source}

The magnetic power to operate the AMR cycle is modeled as the ideal Carnot power, ignoring most thermal irreversibilities:

\begin{equation}
\label{eq:100}
\wmag = \qc \frac{T\ped{H} - \tc}{\tc}
\end{equation}

\noindent while the fluid friction irreversibility is accounted for in the calculation of the pumping power:

\begin{equation}
\label{eq:103}
\wpump = N\ped{r}\frac{1}{\tau}\int_0^{\tau}\frac{\rate{m}\ped{f}}{\rho\ped{f}} \Delta P \diffd{t}
\end{equation}

\noindent where $\Delta P$ is the total pressure drop through the regenerator.

\nomenclature[awm]{$\wpump$}{pumping power [\si{\watt}]}
\nomenclature[awm]{$\wmag$}{magnetic power [\si{\watt}]}
\nomenclature[ap]{$\Delta P$}{pressure drop across one region [\si{\pascal}]}

\subsection{Hydraulic system  and fluid flow profile model}
\label{sec:hydr-syst-model}

The canonical fluid flow profile considered in this work \textcolor{black}{ is the square wave or \emph{instantaneous profile}}, because of the instantaneous change in flow rate, as shown in \autoref{fig:mprofile}. 

\begin{figure}[!ht]
  \centering
  \includegraphics[width=8cm]{mprofile}
  \caption{Instantaneous fluid flow profile}
  \label{fig:mprofile}
\end{figure}

The \textcolor{black}{instantaneous mass flow rate}, \(\rate{m}\ped{f}(t)\), is defined over a cycle with \textcolor{black}{a} period \(\tau\), and represents the fluid flow \textcolor{black}{through} a given regenerator bed. The so-called \emph{hot cycle}, \textcolor{black}{during which} the MCM is magnetized, occupies the \textcolor{black}{time interval} \(0 \le t < \nicefrac{\tau}{2}\), while the \emph{cold cycle}
lies \textcolor{black}{between} \(\nicefrac{\tau}{2} \le t \le \tau\). The flow profile oscillates between two plateaus of equal magnitude
\(\rate{m}\ped{f,max}\) \textcolor{black}{and} opposite directions, \textcolor{black}{which} are centered in each half-cycle. During the hot cycle, the cold
blow period is \(\tau\ped{CB}\), and during the cold cycle the hot blow period is \(\tau\ped{HB}\). If the flow is balanced, then \(\tau\ped{CB} = \tau\ped{HB}\).


During each half cycle, there are periods \textcolor{black}{without}  fluid flow, defined as:

\begin{equation}
\label{eq:48}
\tau\ped{0,HC} = \frac{\nicefrac{\tau}{2} - \tau\ped{CB}}{2}
\end{equation}

\begin{equation}
\label{eq:49}
\tau\ped{0,CC} = \frac{\nicefrac{\tau}{2} - \tau\ped{HB}}{2}
\end{equation}

where HC and CC stand for hot \textcolor{black}{cycle} and cold cycle, respectively.

\nomenclature[vh]{HC}{hot cycle}
\nomenclature[vc]{CC}{cold cycle}

The profile can be mathematically defined as:

\begin{equation}
\rate{m}\ped{f}(t)=
\begin{cases}
0, & 0 \le t < \tau\ped{0,HC} \\
\rate{m}\ped{f,max}, & \tau\ped{0,HC} \le t \le \nicefrac{\tau}{2} - \tau\ped{0,HC} \\
0, & \nicefrac{\tau}{2} - \tau\ped{0,HC} < t < \nicefrac{\tau}{2} + \tau\ped{0,CC} \\
-\rate{m}\ped{f,max}, & \nicefrac{\tau}{2} + \tau\ped{0,CC} \le t \le \tau - \tau\ped{0,CC} \\
0, &  \tau - \tau\ped{0,CC} <  t < \tau \\
\end{cases}
\label{eq:50}
\end{equation}

When the blows have different time durations, the AMR cycle is considered unbalanced, and this \textcolor{black}{is known to have} a negative effect on performance \cite{bib:eriksen16_effec,bib:nakashima18-influen-exp}. In this work, the blows are always balanced, hence the blow fraction, the ratio of blow durations to cycle period \cite{bib:nakashima18-influen-exp}, can be evaluated as:

\begin{equation}
\label{eq:59}
F\ped{B} = \frac{2\tau\ped{B}}{\tau}
\end{equation}

\noindent where $\tau\ped{B}$ is the duration of one blow.

\nomenclature[azt]{$\tau\ped{B}$}{blow duration [\si{\second}]}

The hydraulic system to modulate this fluid flow through different regenerators at different time instants is composed of a pump and a set of electronic valves which can be precisely controlled to yield the desired blow durations. The electrical power consumed by the valve array is computed separately from other work contributions. An application of electronic valves in AMR devices has been presented by \cite{bib:nakashima18-perfor-asses-solen-valves-flow}.

Due to evolving developments in our group, two types of valves are considered, as will be elaborated in \autoref{sec:results-discussions}.

The first approach, called \emph{Type B valves}, uses the  model \textcolor{black}{of} \cite{bib:cardoso16_trans}, assuming that the individual consumption of each valve is independent of frequency and blow fraction. The valve power  $\wvalve$ can be computed as:

\begin{equation}
  \label{eq:58}
  \wvalve = N\ped{v}F\ped{B}\left(\wvalven + \frac{1}{2}\wrelayn\right)
\end{equation}

\noindent where $N\ped{v}$ is the number of valves, $\wvalven$ is the measured average nominal power for one normally-closed electronic valve and $\wrelayn$ is the nominal power for one controlling relay. The factor $\nicefrac{1}{2}$ is due to two valves being controlled by one relay.

\nomenclature[an]{$N\ped{v}$}{number of valves}
\nomenclature[aw]{$\wvalve$}{valve power [\si{\watt}]}
\nomenclature[aw]{$\wvalven$}{nominal power comsumption of one valve [\si{\watt}]}
\nomenclature[aw]{$\wrelayn$}{nominal power comsumption of one relay [\si{\watt}]}

In the second approach, \emph{Type A} valves are used, with lower nominal power but that depends on frequency and blow fraction. For these valves, the valve power was experimentally correlated as:

\begin{equation}
  \label{eq:198}
  \wvalve\,[\si{\watt}]= N\ped{v}\left(0.927 f\,[\si{\hertz}] + 1.023 F\ped{B} + 0.226 f\,[\si{\hertz}] F\ped{B} - 0.037\right)
\end{equation}

Equation~\eqref{eq:198} was correlated for blow fractions of \num{50} and \SI{100}{\percent} and for frequencies in the range \textcolor{black}{of} \num{0.2}--\SI{1.6}{\hertz}, with an uncertainty \textcolor{black}{on} the order of \SI{0.4}{\watt} for a single valve. The use of different valve types will be discussed among the presented results.

Indepedent of the valve type used, it is also assumed that this valve system can produce the fluid flow profile \textcolor{black}{shown in} \autoref{fig:mprofile}, where the displaced fluid mass during one blow in one regenerator bed is $\rate{m}\ped{f,max} F\ped{B} \nicefrac{\tau}{2}$. 

The \emph{utilization factor} can then be calculated as:

\begin{equation}
  \label{eq:110}
  \Phi = \frac{\rate{m}\ped{f,max} F\ped{B} c\ped{s}}{2 f m\ped{s} c\ped{s}}
\end{equation}

\nomenclature[bf]{$\Phi$}{utilization factor}

\textcolor{black}{For an even number of poles in the magnetic circuit}, it may be advantageous to employ an odd number of regenerator beds to  \textcolor{black}{break the symmetry between the magnetic circuit and the beds, thus avoiding positions of equilibrium between magnetic forces which result in torque oscillations} \cite{bib:eriksen}. \textcolor{black}{Additionally}, in general, two valves per regenerator  may be necessary (one at each end), to independently control  the hot and cold blows for each bed and correct flow \textcolor{black}{imbalance}; thus, in all results in this work, the number of valves is \textcolor{black}{calculated as}:

\begin{equation}
  \label{eq:111}
  N\ped{v} = 2 N\ped{r}
\end{equation}

\subsection{Calculation of the coefficient of performance}
\label{sec:calc-coeff-perf}

The coefficient of performance \textcolor{black}{takes} the cooling capacity as \textcolor{black}{the main} output \textcolor{black}{parameter from} the AMR model, \textcolor{black}{in addition to} all previously cited power contributions:

\begin{equation}
  \label{eq:112}
  \cop = \frac{\qc}{\wpump + \wmag + \wvalve}
\end{equation}



\section{Results and Discussions}
\label{sec:results-discussions}


The simulated magnetic refrigerators are composed of multiple prismatic regenerators filled with gadolinium (Gd) spheres as the magnetocaloric solid matrix, with the parameters shown in Table~\ref{tab:params}. The fluid is a mixture of water and ethylene glycol, in the proportion of \SI{80}{\percent}-\SI{20}{\percent} in volume.

\begin{table}[!ht]
  \centering
  \caption{Geometric and operation parameters for the simulated AMR devices}
  \begin{tabular}{c|c}
\hline
    \textbf{Parameter}&\textbf{Value}\\
\hline
Regenerator height & \SI{20}{\mm}\\
Regenerator width & \SI{25}{\mm}\\
Regenerator length & \SI{100}{\mm}\\
Number of regenerators & \num{11} \\
Particle diameter & \SI{0.5}{\mm}\\
Hot source temperature & \SI{298}{\kelvin}\\
Cold source temperature & \SI{278}{\kelvin}\\
\hline
  \end{tabular}
  \label{tab:params}
\end{table}

\subsection{Comparison of the instantaneous and rectified cosine profiles}
\label{sec:comp-inst-cosine}

When comparing the magnetic profiles, the average magnetic field during the hot cycle is the same; this implies in a higher peak for the rectified cosine. For the cold cycle, there are two possible comparison methods:

\begin{enumerate}
\item The minimum values for both profiles is the same;
\item The average value for both profiles during the cold cycle is the same.
\end{enumerate}

As a reference, in all simulations  the minimum value for the rectified cosine was fixed at $\average{B}\ped{min,RC} = \SI{0.1}{\tesla}$. For the first comparison, the minimum value for the instantaneous value was kept at the same level ($\average{B}\ped{min,I} =  \average{B}\ped{min,RC}$), hence, the average low value for the RC profile was higher ($B\ped{l,RC} > B\ped{l,I}$). Although this does not seem a fair comparison, in practice it is feasible to generate a profile with a constant low-valued plateau for the magnetic field \cite{bib:insinga16_optim,bib:benedict16_desig}. For the other comparison method, the low level of the instantaneous profile was increased ($\average{B}\ped{min,I} >  \average{B}\ped{min,RC}$) to keep the average the same ($B\ped{l,RC} = B\ped{l,I}$). In addition, simulations were ran for various values of blow fraction and the optimal cases were selected; the rectified cosine profile can benefit from a smaller blow fraction because this can concentrate the flow during the periods of very high or low fields. In this section, all results show the resulting optimal case of the blow fraction for each profile (for the values tested): fluid flowing during the entire period for the instantaneous profile, and only during \SI{60}{\percent} of the period for the cosine profile.


Figure~\ref{fig:cos_ins} shows the cooling capacity attained by the device at a frequency of \SI{1}{\hertz} and different utilizations, for both comparison methods.  The horizontal axis shows the average value during the high field region. Notice how the results are conflicting; for Fig.~\ref{fig:comp_min}, the instantaneous profile almost always yields a higher performance, while for Fig.~\ref{fig:comp_avg} the rectified cosine profile generates higher cooling capacities. Since the average during the hot cycle (high field stage) is the same, the main difference is due to the low magnetic field levels. For the ``same minimum'' comparison, the instantaneous profile is capable of keeping the magnetic field at low levels for the whole half-cycle, which is beneficial for performance;  for $\Phi=1.0$ and $B\ped{h} = \SI{1.40}{\tesla}$, the cooling capacity for the instantaneous profile is $\SI{196.3}{\percent}$ higher than for the cosine profile.  As demonstrated by \cite{bib:asme-mce}, a higher average magnetic field during the low-field stage increases the solid temperature and consequently results in  warmer fluid entering the cold heat exchanger, representing a thermal loss. In the ``same-average'' comparison, the cosine profile is capable of achieving much lower levels (since its minimum value is fixed), therefore providing more cooling power. However, as explained before, it is possible to design near-instantaneous profiles with very low levels of the magnetic field.


\begin{figure}[!ht]
  \centering
%  \subfloat[Same minimum]{\includegraphics[width=8cm]{Qc_B_comp_f_1_same_minimum}\label{fig:comp_min}}
\,
%  \subfloat[Same average]{\includegraphics[width=8cm]{Qc_B_comp_f_1_same_low_average}\label{fig:comp_avg}}
\,
  \caption{Cooling capacity as a function of the  average high magnetic field, for different utilizations. “I”: instantaneous (blow fraction of \SI{100}{\percent}); “RC”: rectified cosine  (blow fraction of \SI{60}{\percent}), for both comparison methods.}
 \label{fig:cos_ins}
\end{figure}

The only exception in the ``same minimum'' comparison is seen for the lowest utilization of $\Phi = 0.2$, where the performance is slightly better for the RC profile. Since the blow fraction for the cosine is lower, the mass flow rate must be higher to yield the same utilization (cf. Eq.~\eqref{eq:2}), and this increases heat transfer rate, as previously explained ---  outweighing the effects of the magnetic field.

In general, the performance of the AMRs operating with the instantaneous magnetic profile yields better results. As can be seen on Fig.~\ref{fig:cos_ins}, an instantaneous profile with the lowest possible value of $\bmin$ and the highest possible value of $\bmax$, with a flow profile occupying the whole cycle with average values of utilization, results in the highest values of cooling capacity among all simulations, therefore, representing the ideal magnetic profile for an AMR.

\subsection{Analysis of the instantaneous profile}
\label{sec:deta-analys-inst}

As shown in the previous section, the instantaneous profile yields the highest values for cooling capacity. Therefore, in this section, a more detailed analysis of such profile has been carried out. Figure~\ref{fig:qc_phi_inst} shows the cooling capacity as function of utilization, for several levels of the maximum magnetic field for the instantaneous profile, for two different operating frequencies. Because of the conflicts between a low heat transfer rate for flow rates too low and a loss in regenerator effectiveness in flow rates too high, there are critical values of utilization that maximize cooling capacity, and this critical value grows with the level of magnetic field. With higher magnetic field, the increase in the MCE surpasses the loss of effectiveness, and one can go to higher flow rates without losing performance. It can also be seen in Fig.~\ref{fig:qc_phi_inst} that at higher frequencies, the values of cooling power are higher, and also the critical values of utilization are lower; however, this usually is achieved at the cost of higher power consumption in AMR devices at higher frequencies \cite{bib:lei15_study,NIKNIA2016601}.

\begin{figure}[!ht]
  \centering
\subfloat[$f = \SI{1}{\hertz}$]{\includegraphics[width=8cm]{Qc_Phi_inst_f_1_Hmin_005_FB_100}\label{fig:Qc_phi_inst_1}}
\,
  \subfloat[$f = \SI{2}{\hertz}$]{\includegraphics[width=8cm]{Qc_Phi_inst_f_2_Hmin_005_FB_100}\label{fig:Qc_phi_inst_2}}
  \caption{Cooling capacity as a function of utilization, for various values of the high magnetic field  of the instantaneous profile}
  \label{fig:qc_phi_inst}
\end{figure}

\begin{figure}[!ht]
  \centering
  \includegraphics[width=8cm]{Qc_H_inst_f_1_Phi_40}
  \caption{Cooling capacity as a function of  regenerator height (all other parameters from Table~\ref{tab:params} constant) for various values of the high magnetic field of the instantaneous profile}
\label{fig:Qc_H_inst}
\end{figure}

The parameters in Table~\ref{tab:params} were chosen from preliminary simulations, but they are not guaranteed to be the optimal choices. To understand the impact of regenerator geometry on the performance with the instantaneous profile, the regenerator height was varied in Fig.~\ref{fig:Qc_H_inst}, and all other parameters from Table.~\ref{tab:params} were kept fixed. As expected, higher magnetic fields allow for smaller regenerators (and hence more compact systems). For instance, to achieve a capacity of \SI{100}{\watt}, increasing the field from \num{1.0} to \SI{1.2}{\tesla} result in using regenerators \SI{36}{\percent} smaller.

\section{CONCLUSIONS}
\label{sec:conclusions}

A one-dimensional AMR numerical model was used to compare the performance of a magnetic refrigerator under different operating and geometric parameters and with the instantaneous and rectified cosine magnetic profiles. Most of the AMRs operating with an instantaneous magnetic profile have resulted in better performances than those with a cosine profile, except in cases of low utilization. Although the cosine profile reaches higher levels of the magnetic field, the instantaneous profile can keep the magnetization (and demagnetization) for a longer period, which is shown to be beneficial to performance. Further analysis of the instantaneous profile has shown that the use of the highest possible value of the maximum magnetic field allows for using smaller regenerators and lower utilization, without losing performance.

\begin{acknowledgements}
This work was supported by CNPq, Embraco and EMBRAPII Unit Polo/UFSC.  
\end{acknowledgements}

% BibTeX users please use one of
%\bibliographystyle{spbasic}      % basic style, author-year citations
%\bibliographystyle{spmpsci}      % mathematics and physical sciences
\bibliographystyle{spphys}       % APS-like style for physics

\bibliography{thermo-ref/Thermo-Foam-Ref,thermo-ref/thesis}

\end{document}

%%% Local Variables:
%%% mode: latex
%%% TeX-master: t
%%% End:
